
% =====================================================================
\section{Predictions}
\label{sec:predictions}
% =====================================================================

The model has three parameters, the vacuum value $\vVS$, the Yukawa coupling $\yVS$ and the
self-coupling $\lamVS$ of the Vorton-Scalar or, equivalently, the fermion mass $m$, the mass $m_s$
of the radial mode and the coupling $\gVS=m/\vVS$. Everything else is derived. This section collects
the quantitative consequences that the preceding sections derived, so that a reader who wants to
test the model, or to extend it, finds them in one place. Table~\ref{tab:predictions} lists them
with the figure or table that displays each. They are of three kinds: predictions for the matter
that the model describes, cosmological predictions, and predictions that distinguish the two
spacetimes.

\begin{table*}
\caption{\label{tab:predictions}Quantitative predictions of the model. Unless a spacetime is named, a
prediction holds for both. $C=\gVS^2m^2/m_s^2$ and $\DVS=\aVS m/m_s$.}
\begin{ruledtabular}
\begin{tabular}{>{\raggedright\arraybackslash}p{0.035\textwidth}>{\raggedright\arraybackslash}p{0.33\textwidth}>{\raggedright\arraybackslash}p{0.37\textwidth}>{\raggedright\arraybackslash}p{0.18\textwidth}}
 & Quantity & Prediction & Displayed in \\
\hline
\multicolumn{4}{l}{\emph{Matter}}\\
1 & masses and couplings & $m=\yVS\vVS/\sqrt2$, $m_s=\sqrt{2\lamVS}\,\vVS$, massless phase mode, $\gVS=m/\vVS$, $\GVS/\sqrt2=1/2\vVS^2$ & Eqs.~(\ref{eq:masses}), (\ref{eq:LNC}) \\
2 & force of the neutral current between two fermions & contact attraction $-(\GVS/2\sqrt2)\,m^2/E^2$ in the vector channels along the extra axes; repulsion of the same strength in the two-form channels & Fig.~\ref{fig:kernels}, Table~\ref{tab:kernels} \\
3 & channel of the largest gap & the vector channel; at weak coupling its gap exceeds that of the two-form channel by a factor {{GEN:gap_ratio_lo}} to {{GEN:gap_ratio_hi}} & Fig.~\ref{fig:gap}, Table~\ref{tab:gap} \\
4 & bound pair of two fermions & exists if and only if $\DVS>{{GENM:dimer_dc}}$ & Fig.~\ref{fig:dimer}, Table~\ref{tab:threshold} \\
5 & thermodynamics of the paired state & $\Delta(0)/T_c={{GENM:cjt_gap0_over_tc}}$; jump of the heat capacity ${{GENM:cjt_c_jump}}\,C_n$ & Fig.~\ref{fig:thermo}, Table~\ref{tab:thermo} \\
6 & gaps of a general state & two fourfold gaps, Eq.~(\ref{eq:twogaps}); ground states with one eightfold gap; a stationary state with four gapless internal states and half the condensation energy & Fig.~\ref{fig:landscape}, Table~\ref{tab:landscape} \\
7 & collective modes & four gapless modes: the sound, with $c_s^2=v_F^2/3$ at weak pairing, and three rotations of the order parameter & Table~\ref{tab:landscape}, Eq.~(\ref{eq:sound}) \\
8 & fermion mass in matter & $m^*=m(1-Cn/m^3)$ at low density; $m^*=0$ above $\kF=\pi m/2\sqrt C$ (weak coupling) & Fig.~\ref{fig:normal}, Table~\ref{tab:normal} \\
9 & self-bound matter & if $C>3\pi^2/8={{GENM:cjt_c_crit}}$, or $3\pi^2/5={{GENM:cjt_c_crit_sea}}$ with the negative-energy states & Fig.~\ref{fig:normal}, Table~\ref{tab:normal} \\
10 & stability of the vacuum at one loop & lost for $C\ge3.7$: the potential falls below its vacuum value beyond {{GEN:cjt_vacstab_lo}} to {{GEN:cjt_vacstab_hi}} $\vVS$ & Fig.~\ref{fig:vacpot}, Table~\ref{tab:vacstab} \\
11 & melting of the fermion mass & at $T_{1/2}$ between {{GEN:cjt_thermal_thalf_m_lo}} and {{GEN:cjt_thermal_thalf_m_hi}} $m$; estimates (\ref{eq:Tc}) & Fig.~\ref{fig:thermal}, Table~\ref{tab:thermal} \\
12 & pair condensate in the vacuum & only if $\aVS>\alpha_c$, with $\alpha_c$ between {{GEN:cjt_alpha_c_min}} and {{GEN:cjt_alpha_c_max}}, and $\pi/2$ for an infinite cutoff & Fig.~\ref{fig:vacuum}, Table~\ref{tab:vacuum} \\
\multicolumn{4}{l}{\emph{Cosmology}}\\
13 & step of the dark radiation & factor $15^{1/3}=2.466$ in the density when the fermions annihilate & Fig.~\ref{fig:step}, Eq.~(\ref{eq:rstep}) \\
14 & cost of the local Hubble constant & $\sqrt{\Delta\chi^2}\ge$ {{GEN:dr_sig_local_min_all}}; {{GEN:dr_sig_local_min_model}} to {{GEN:dr_sig_local_max_model}} for a step in the epochs probed; gain at $\Delta\chi^2=4$ at most {{GEN:dr_gain_2sigma_best}} $\kmsMpc$ & Figs.~\ref{fig:dr}, \ref{fig:zt}, Table~\ref{tab:dr} \\
15 & helium with an early step & mass fraction {{GEN:dr_yhe_early}} instead of {{GEN:dr_yhe_nostep}} at $N=0.5$ & Sec.~\ref{sec:hubble} \\
16 & radiation of the phase mode & $\Delta N_{\rm eff}$ between {{GEN:dr_weinberg_lo}} and {{GEN:dr_weinberg_hi}}, by the epoch of decoupling & Eq.~(\ref{eq:weinberg}) \\
17 & small-scale structure with superfluid dark matter & suppressed; requires $m>{{GENM:sfdm_mmin_5p3}}$ keV, or a condensate with $R_{TF}<{{GENM:sfdm_rtf_pc_5p3}}$ pc & Fig.~\ref{fig:sfdm}, Tables~\ref{tab:sfdm}, \ref{tab:sfdmbounds} \\
18 & superfluidity in halos & absent: degeneracy needs $m<{{GENM:dm_mdeg_dwarf}}$ keV, condensation of pairs $2m<{{GENM:dm_mbmax_dwarf}}$ eV, in a dwarf galaxy & Fig.~\ref{fig:dmwindows}, Table~\ref{tab:halos} \\
19 & bound pairs as dark matter & $m>{{GENM:dm_mmin_alpha1}}$ MeV for $\aVS=1$, from the limit on the self-interaction & Table~\ref{tab:selfint} \\
\multicolumn{4}{l}{\emph{The two spacetimes}}\\
%%<1644>
20 & dependence on the additional times ($4+4$) & modes with $\lvert\kappa\rvert^2>m^2+\lvert k\rvert^2$ grow as $\exp\big(\sqrt{\lvert\kappa\rvert^2-m^2-\lvert k\rvert^2}\;x4\big)$ & Eq.~(\ref{eq:K2}) \\
21 & relative phase of the components of the order parameter & $\pm90^\circ$ between the hidden axis and a time-like axis ($4+4$); $0$ ($7+1$) & Fig.~\ref{fig:landscape} \\
22 & the three rotational gapless modes & massive as soon as the fields depend on $x0$ ($4+4$); exact Goldstone modes ($7+1$) & Sec.~\ref{sec:ground} \\
23 & electroweak pattern & singlet of the rotations of the additional times; $(2,2)$ only under a hidden symmetry ($4+4$) & Table~\ref{tab:verdict} \\
%%</1644>
%%<1671>
20 & dependence on the extra coordinates ($7+1$) & a momentum $\vec p_\perp$ along them adds to the mass, $E^2=m^2+\lvert\vec p\rvert^2+\lvert\vec p_\perp\rvert^2$; on the $4+4$ spacetime a dependence on the additional times produces growing modes & Eq.~(\ref{eq:hamiltonian}) \\
21 & relative phase of the components of the order parameter & $0$ ($7+1$); $\pm90^\circ$ between the hidden axis and a time-like axis ($4+4$) & Fig.~\ref{fig:landscape} \\
22 & the three rotational gapless modes & exact Goldstone modes ($7+1$); massive as soon as the fields depend on $x0$ ($4+4$) & Sec.~\ref{sec:ground} \\
23 & electroweak pattern & a $(2,2)$ of the rotations of the extra axes with the charges $(+1,0,0,-1)$, and the content of two Higgs doublets ($7+1$) & Table~\ref{tab:verdict} \\
%%</1671>
\end{tabular}
\end{ruledtabular}
\end{table*}

\subsection{Predictions for matter}

Rows 1 to 12 describe a kind of matter: fermions with eight internal states whose mass comes from a
scalar field. Their status differs. Rows 1, 2, 4 and 6 are exact consequences of the Lagrangian
within the stated order: the relations between masses and couplings, the amplitudes of one-boson
exchange, the threshold of the two-body problem, and the algebra of the gap matrix. Rows 3 and 5 hold
in the weak-coupling regime, in which the gap is small compared with the Fermi energy; row 5 is the
universal behavior of any weakly coupled superfluid with a single gap, and it tests the pairing
mechanism only through the statement that the gap \emph{is} single, that is, that the ground state is
unitary. The characteristic predictions are rows 2, 3 and 6: a neutral current that attracts in the
vector channels and repels in the two-form channels, so that the order parameter is a four-vector
and not an antisymmetric tensor; and a spectrum with two fourfold gaps away from the ground state.
An experiment on such matter that could resolve the eight internal states would see the half-paired
state as a state in which four of them remain gapless.

Rows 8 to 12 concern the scalar field in a medium, at finite temperature and in the vacuum. They are
mean-field results. Rows 8 and 9 depend on one number, $C$, and they imply a qualitative prediction
that does not depend on details: dense matter of these fermions is in the abnormal state, in which
they are massless, and for strong coupling it is bound without gravity. Row 10 limits the range in
which rows 9 and 12 can be trusted.

Two further consequences are taken from the literature and were not computed here
(Sec.~\ref{sec:vortons}): the vortices of the superfluid carry the circulation $\pi/m$ per unit of
winding, and the vortex lines of the Vorton-Scalar bind fermions of zero mass that can carry a
persistent current.

\subsection{Cosmological predictions}

Rows 13 to 19 were obtained by simulation and are compared with observation in
Secs.~\ref{sec:hubble} and \ref{sec:dm}. They are the predictions on which hypothesis H3 fails. Row
13 is a prediction in the strict sense, a pure number that follows from the count of the degrees of
freedom, sixteen fermionic and one bosonic. It distinguishes the model from the stepped fluid of
Ref.~\cite{Aloni2022}, whose step is $(15/7)^{1/3}$, and it is the reason why the model fares worse
than that fluid when the step falls into the epochs that the microwave background probes. Row 15 is
a prediction that could be tested independently of the microwave background, by the primordial
helium abundance, if the dark radiation had the density needed; the difference is
{{GEN:dr_yhe_diff}} in the mass fraction at $N=0.5$, which is comparable to the uncertainty of about
0.003 of present determinations \cite{PDG2024}. Rows 17 to 19 state, in
numbers, why the superfluid cannot act on galactic scales.

\subsection{Predictions that distinguish the two spacetimes}

Rows 20 to 23 are the places where the signature of the eight-dimensional metric reaches observable
physics. In the sector of fields that depend on ordinary space and time only, the two models have
the same internal symmetry algebra (Table~\ref{tab:symmetry}), the same exchange amplitudes
(Table~\ref{tab:kernels}) and the same thermodynamics; a measurement within that sector that is
insensitive to the relative phases of the components of the order parameter cannot tell them
apart. They differ in three respects. The metric on the four extra axes enters the energy of the
superfluid through the quadratic form $\eta_{ab}\Delta^a\Delta^b$ of Eq.~(\ref{eq:quartic}), which
fixes the relative phase of the components in the ground state (row 21). The symmetry that rotates
the hidden axis into the other three is a symmetry of the spacetime on $7+1$ and a property of the
zero-mode sector on $4+4$, which decides whether the three rotational modes are exact Goldstone
modes (row 22) and whether the electroweak pattern can have a geometric origin (row 23). And the
dependence of the fields on the three additional coordinates is harmless on $7+1$, where it adds to
the mass, and destroys the quantum theory on $4+4$, where it produces growing modes (row 20).
